\chapter{Matrix strategies}

Matrix jobs run the same job across multiple configurations.
They are ideal for OS, language versions, or test shards.

\section{A basic matrix}
\begin{lstlisting}[language=YAML]
jobs:
  test:
    runs-on: ${{ matrix.os }}
    strategy:
      matrix:
        os: [ubuntu-latest, macos-latest, windows-latest]
        node: ["18", "20"]
    steps:
      - uses: actions/checkout@v4
      - uses: actions/setup-node@v4
        with:
          node-version: ${{ matrix.node }}
      - run: npm test
\end{lstlisting}

\section{Include and exclude}
Use \code{include} for one-off combinations and \code{exclude} for invalid ones.

\begin{lstlisting}[language=YAML]
strategy:
  matrix:
    os: [ubuntu-latest, windows-latest]
    python: ["3.11", "3.12"]
    include:
      - os: ubuntu-latest
        python: "3.12"
        experimental: true
    exclude:
      - os: windows-latest
        python: "3.12"
\end{lstlisting}

\section{Allowing experimental failures}
Mark experimental runs as optional with \code{continue-on-error}.

\begin{lstlisting}[language=YAML]
- name: Run tests
  continue-on-error: ${{ matrix.experimental || false }}
  run: ./scripts/test
\end{lstlisting}

\section{Controlling parallelism}
Large matrices can overwhelm runners.
Use \code{max-parallel} to limit concurrency.

\begin{lstlisting}[language=YAML]
strategy:
  fail-fast: false
  max-parallel: 4
  matrix:
    shard: [1, 2, 3, 4, 5, 6, 7, 8]
\end{lstlisting}

\section{Dynamic matrices}
You can generate a matrix from a previous step and parse it with \code{fromJSON}.
This keeps the matrix small and relevant.

\begin{lstlisting}[language=YAML]
- id: plan
  run: echo 'matrix={"service":["api","web"]}' >> "$GITHUB_OUTPUT"

strategy:
  matrix: ${{ fromJSON(steps.plan.outputs.matrix) }}
\end{lstlisting}

\section{Matrix design heuristics}
A matrix should represent real product dimensions.
Avoid adding variants just because they are easy to add.

\begin{itemize}[nosep]
  \item Use matrices for compatibility validation.
  \item Use shards for speed, not correctness.
  \item Keep a default job for the most common case.
\end{itemize}

\section{Fail-fast tradeoffs}
Fail-fast is good for quick feedback but hides the full failure surface.
Disable it when you need complete signal across all variants.

\section{Artifacts per variant}
If each variant builds a binary, upload artifacts with the matrix values in the name.
This avoids collisions and makes debugging easier.

\begin{checklistbox}{Checklist: matrix}
\begin{itemize}[nosep]
  \item Use matrices for the smallest set of meaningful variants.
  \item Avoid multiplying dimensions without a strong reason.
  \item Mark experimental variants explicitly.
\end{itemize}
\end{checklistbox}
